%Keywords: Newton polyhedron, algebraic equations, differential equations,
%asymptotics, first approximation, singular perturbation, boundary layer.

\documentstyle[12pt]{article}
\setlength{\evensidemargin}{0in}
\setlength{\oddsidemargin}{0in}
\setlength{\textwidth}{16.cm}
\setlength{\textheight}{8.5in}
\setlength{\topmargin}{0in}
\setlength{\headheight}{0in}
\setlength{\headsep}{0in}
\setlength{\itemsep}{-\parsep}
\renewcommand{\topfraction}{.9}
\renewcommand{\textfraction}{.1}
\newcommand{\ol}{\setlength{\itemsep}{0pt.}\begin{enumerate}}
\newcommand{\eol}{\end{enumerete}\setlength{\itemsep}{-\parsep}}
\newcommand{\R}{{\rm{I\! R}}}
\newcommand{\Q}{{\rm Q\!\!\! l}}
\newcommand{\C}{{\rm{I\!\!\! C}}}
\newcommand{\N}{{\rm{I\! N}}}
\newcommand{\defi}{{\;\buildrel \rm def \over =}\;}
\hyphenation{or-di-na-ry}
\hyphenation{ap-pro-x-i-ma-t-i-o-n}
\hyphenation{per-tur-ba-ti-on}
\hyphenation{si-n-g-u-la-ri-ti-es}
\hyphenation{di-f-fe-re-n-ti-al}

\begin{document}
\Large
\centerline{\bf The Newton polyhedron in the nonlinear Analysis}

\centerline{\bf Alexander D. Bruno}
\centerline{Institute of Applied Mathematics, Moscow}
\centerline{125047 Russia (bruno@applmat.msk.su)}
\bigskip
\bigskip

\centerline{\bf Abstract}

We give a simple presentation of an algorithm of selecting asymptotical
first approximations of equations (algebraic and ordinary differential
and partial differential). Here the first approximation of a solution
of the initial equation is a solution of the corresponding first
approximation of the  equation. The algorithm
is based on the geometry of power exponents including the Newton
polyhedron. We give also a surway of applications of the algorithm in
problems of Celestial Mechanics and Hydrodynamics.

{\bf Keywords:} Newton polyhedron, algebraic equations, differential equations,
asymptotics, first approximation, singular perturbation, boundary layer.

\bigskip

\centerline{\bf Introduction}
\bigskip

Let $F(X)$ be a polynomial (may be vectorial) of some variables $x_1, ..., x_l$
and their derivaties (ordinary or partial). Let $X=G(T)$ be a solution of the
equation $F(X)=0$, where $T=(t_1,\dots, t_m)$ are some parameters. We are
interested in solutions for which some of $x_i\to 0$ or $\infty$. And we want
to find an asymptotics of such a solution $G(T)$.
Let $\hat G(T)$ be the asymptotical first approximation of the solution
$G(T)$ and let $\hat F(X)$ be the corresponding first approximation of the
polynomial $F(X)$.

{\bf General Statement:} {\it The first approximation $X=\hat G(T)$ of the
solution $X=G(T)$ is a solution of the corresponding first approximation
$\hat F(X)=0$ of the equation $F(X)=0$.}

The aim of the talk is to explain the algorithm for selecting first
approximations, to show how this approach works for algebraic equations and
for differential equations, and to point out some applications.
\bigskip

\centerline {\bf 2.~A polyhedron and normal cones of its faces}
\bigskip

Let in the space $\R^{n}$ with Cartesian coordinates
$q_{1}, \dots , q_{n}$ we have a finite set $D$ of points
$Q_{j}=(q_{1j},\dots, q_{nj}),\: j=1, \dots, m$. Let ${\R}_*^{n}$
denote the dual space such  that for $P=(p_{1},\dots, p_{n}) \in
{\R}_*^{n}$ and $Q=(q_{1},\dots, q_{n}) \in \R^{n}$ there is the scalar
product $\langle{P,\, Q}\rangle = p_{1}q_{1}+\dots+p_{n}q_{n}$. For a fixed
vector $P\neq 0$ let $D_{P}$ denote such subset of the set $D$ on which
the scalar product $\langle{P,\,Q}\rangle$ has the maximal value $c$, that is
$$
\begin{array}{l}
\langle{P,\, Q_{j}}\rangle = c \;\;{\rm for}\;\; Q_{j} \in
D_{P}\;\; {\rm and}\\
\langle{P,\, Q_{j}}\rangle < c \;\;{\rm
for}\;\; Q_{j} \in D\setminus D_{P}.
\end{array}\eqno(1)
$$

Now for the given set $D$ and for each $P\neq 0$ we will find the
corresponding set $D_P$.
Let $M$ denote the convex null of the set $D$. The boundary $\partial M$ of
the {\it polyhedron} $M$ consists of {\it faces} $\Gamma_k^{(d)}$ of
different dimension $d$ $(0\leq d<n)$.  Zerodimensional faces
$\Gamma_k^{(0)}$ are verteces of the polyhedron $M$, onedimensional faces are
its edges etc.
Accordingly (1) to each vector $P\in \R_*^n$ there corresponds in $\R^n$ a
{\it supporting} to the set $ D$ {\it hyperplane} $L_P=\{Q:\langle
P,Q\rangle=c\}$.
It intersects the polyhedron $M$ along a face $\Gamma_k^{(d)}$. The set
$U_k^{(d)}$  of all vectors $P\in \R^n_*$ such, that $L_P\cap
M=\Gamma_k^{(d)}$, is called the {\it normal cone} of the face
$\Gamma_k^{(d)}$.
The normal cone $U_k^{(n-1)}$ of the hyperface $\Gamma_k^{(n-1)}$ is a ray
ortogonal to the face and directed out of the polyhedron $M$. The normal cone
$U_k^{(n-2)}$ is a sector bounded by rays $U_l^{(n-1)}$   and $U_m^{(n-1)}$
if $\Gamma_k^{(n-2)}=\Gamma_l^{(n-1)}\cap\Gamma_m^{(n-1)}$ etc. The union of
normal cones of all faces is $\R^n_*\setminus\{0\}$. Let us consider the {\it
boundary subsets} $D_k^{(d)}=\Gamma_k^{(d)}\cap D$ of the set $D$.

{\bf Theorem 1 [1].} {\it If $P\in U_k^{(d)}$, then $D_P=D_k^{(d)}$.}

{\bf Example 1.} For $n=2$ and
$$
D=\{(3,0), (0,3), (1,1)\}\eqno(2)
$$
the polygon  $M$, faces $\Gamma_k^{(d)}$ and their normal cones are shown in
Figures 6 and 7 of the book [1]. Here
$$\begin{array}{lll}
D_1^{(0)}=Q_1,&D_2^{(0)}=Q_3,&D_3^{(0)}=Q_2;\\
D_1^{(1)}=\{Q_1,Q_3\},&D_2^{(1)}=\{Q_2,Q_3\},&D_3^{(1)}=\{Q_1,Q_2\};\\
U_1^{(1)}=-\lambda(1,2),&U_2^{(1)}=-\lambda(2,1),&U_3^{(1)}=-\lambda(1,1)\;\; 
(\lambda>0).
\end{array}
$$
\bigskip

\centerline{\bf 3.~Algebraic equations}
\bigskip

Let us consider a polynomial
$$
f(X) =  \sum\limits_{j=1}^{m} f_{j}X^{Q_j},
$$
where $X^{Q}=x_{1}^{q_1}\dots x_{n}^{q_n}$ and $D=\{Q_{1},\dots,Q_{m}\}$.
The convex hull $M$ the set $D$ is called the {\it Newton polyhedron} of the
polynomial $f$, and we can construct all accompanying objects as described
abowe.
Let $D_k^{(d)}$ is a boundary subset of the set $D$.
The sum
$$
\hat f_{k}^{(d)}(X) = \sum f_{j}X^{Q_j}\;\;{\rm for}\; Q_{j} \in
D_{k}^{(d)} \eqno(3)$$
is called the {\it truncation} of the polynomial $f(X)$.
Now we consider a curve of the form
$$
x_{i} = b_{i}\tau^{p_i}(1 + o(1)), \;\; b_i\neq 0, \;\;i =
1,\dots,n,\eqno(4) $$
where $\tau \to \infty$. On such curve a monomial
$$
X^{Q} = B^{Q}\tau^{\langle P,Q \rangle}(1+o(1)),
$$
where $B=(b_{1},\dots,b_{n})$ and $P=(p_{1},\dots,p_{n})$. If   $P\in
U_{k}^{(d)}$, then on the curve (4)
$$
f(X) = \hat f_{k}^{(d)} (B)\tau ^{c} + \tau ^{c}o(1),
$$
where $c$ is from (1). Thus, the truncation (3) is the first approximation of
the polynomial  $f(X)$ on curves (4) with  $P\in U_k^{(d)}$. If curve (4) is
a solution of the equation $f(X)=0$ then
$\hat f_{k}^{(d)}(B) = 0$. That is we have

{\bf Theorem 2} [1,2,4]. {\it The first approximation
$x_{i} = b_{i} \tau^ {p_i}, \;\; i=1,\dots, n$ of the solution (4) is a
solution of the corresponding truncated equation
 $\hat f_{k}^{(d)}(X)=0$.}

This approach allows to  find solutions of the equation $f=0$ with any degree
of accuracity [1,2,4].

{\bf Example 2.} Let
$f=x_1^3+x_2^3-2x_1x_2$, then the set $D$ is (2), and among truncations
$\hat f_{k}^{(d)}$ there are
$$
\begin{array}{c}
\hat f_{1}^{(0)}=x_1^3,\quad
\hat f_{2}^{(0)}=-2x_1x_2,\quad
\hat f_{3}^{(0)}=x_2^3;\\
\hat f_{1}^{(1)}=x_1^3-2x_1x_2,
\hat f_{2}^{(1)}=x_2^3-2x_1x_2.
\end{array}
$$
If (4) is a solution of the equation $f=0$ with $P\neq 0$, then $P$ cannot ly
in $U_k^{(0)}$ because $B^Q\neq 0$. If on the solution (4)
 $x_1,x_2\to 0$ for
$\tau\to\infty$, then $P\in U_1^{(1)}$ or $U_2^{(1)}$.  If
$P\in U_1^{(1)}$, i.e. $2p_1=p_2<0$,
then the first approximations $x_i=b_i\tau^{p_i}$ $(i=1,2)$ satisfies
the truncated equation  $x_1^3-2x_1x_2=0$, i.e $x_2=x^2_1/2$.
Similarly, if $P\in U_2^{(1)}$,
then the first approximation of the solution (4) is $x_1=x_2^2/2$. That is
two brauches of the curve $f=0$ go trough the zero $x_1=x_2=0$
(see [1], Figure 19).
Other applications of the Newton polyhedron for algebraic equations
see in [1,2,4,20]
\bigskip

\centerline{\bf 4.~Ordinary differential equations}
\bigskip

Now we consider a differential polynomial
$$
g(z_{1},\dots,z_{n}) =\sum g_{R}Z^{R},\eqno(5)
$$
where $Z^{R} = z_{1}^{r_1}\dots z_{n}^{r_n}$ and
$$
z_{1} = x, \;\;z_2=y,\;\; z_{2+l} = d^{l}y/dx^{l}, \;\;
l=1,\dots,n-2.\eqno(6)
$$
To calculate powers of variables $x$ and $y$ we shall consider
$d^{l}y/dx^{l}$ as the ratio $y/x^l$. Hence, to each monomial $X^R$ we put in
correspondance the vector power exponent
$Q=(q_{1},q_{2})=Q(R)$ with  coordinates
$$
q_{1} = r_{1} - \sum_{l=1}\limits^{n-2} lr_{2+l},\;\; q_{2}=
r_{2}+\dots+r_{n},\eqno(7)
$$
i.e. the point $Q$ of the plane $\R^2$. After the substitution
$$
x= z_1= \tau^{p_1},\;\; y=z_2=b\tau ^{p_2}
$$
we have  $Z^{R} = {\rm const}\cdot \tau ^{\langle P,Q \rangle}$, where
$P=(p_1,p_2)$.
Now the set $D=D(g)$ is the set of all points $Q(R)$ with $g_{R}\neq 0$.
Using boundary subsets $D_k^{(d)}$ of the set $D$, we define {\it
truncations
$$
 \hat g_{k}^{(d)}(Z) = \sum g_{R}Z^{R}\;\,\,\,{\rm for}\,\,\; R: Q(R)
\subset D_{k}^{(d)}
$$
of the  differential polynomial (4).}

{\bf Theorem 3 }[3]. {\it Let the differential equation $g(Z)=0$ have a
solution of the form
$$
x_{1}=\tau ^{p_1},\; y = \tau ^{p_2}(b+\sum_{i=1}\limits ^{\infty}
b_{i} \tau ^{s_i}),\quad b\neq 0,\eqno(9)
$$
where
 $P=(p_1,p_2)\neq 0$, $\tau \rightarrow \infty$   and
$0> s_i> s_{i+1}$.
If  $P \in U_{k}^{(d)}$ ,  then the first approximation (8) of the solution
(9)  is a solution of the corresponding truncated equation
 $\hat f_{k}^{(d)}(X)=0$.
 }

 Examples and applications of the Newton polyhedron in ordinary differential
 equations (including singular perturbations) see in [1, 3, 5--10, 13--19,
 21--25].

\bigskip
\centerline{\bf 5.~Partial differential equations}
\bigskip

Let us consider the partial differntial  polynomial
$$
h(W)=\sum h_R W^R,
\eqno(10)
$$
where $W=(w_1,\dots, w_n)$, $w_1=x$, $w_2=y$ and $w_{j+2}$
are partial derivations of the form
 $(\partial^{k+l}\psi)/(\partial x^k\partial y^l)$
To calculate the power exponents, we consider such a derivative as the
monomial
$x^{-k}y^{-l}\psi$, i.e. to it there corresponds the vector exponent
$Q=(q_1,q_2,q_3) = (-k,-l,1)\in \R^3$.
To a product of such derivatives there corresponds the sum of such vector
exponents. Thus, to each monomial $W^R$ there corresponds a point
$Q=Q(R)\in \R^3$, and to the polynomial (10) there corresponds a set $D$
of such points $Q$. To each boundary subset
$D_k^{(d)}$ there corresponds the {\it truncated polynomial}
$$
\hat h_k^{(d)} =\sum h_RW^R\quad {\rm for}\quad R:Q(R)\in D_k^{(d)}.
$$
It is the first approximation of the polynomial (10) when vector
$(\log|x|, \log|y|, \log|\psi|)$ tends to infinity near the normal cone
$U_{k}^{(d)}\subset \R^3_*$.

For an expansion
$\psi = f(x,y)$ and for a given vector $(p_1,p_2)$ we can define its first
approximation
$\hat f_{(p_1,p_2)}$ as in Section 3. On the curves (4) it has the order of
magnitude
const$\cdot\tau^{p_3}$.  To the vector  $P=(p_1,p_2,p_3)$ there corresponds
the truncation
$\hat h_P(W)$ of the differential polynomial  (10).

{\bf Theorem 4} [3]. {\it If
 $\psi=f(x,y)$ -- is a solution of the equation $h(W)=0$, then its first
 approximation
$\psi=\hat f_{(p_1,p_2)}$ is a solution of the corresponding truncated
equation $\hat h_P(W) = 0$.}

 {\bf Example 3 } [11]. The Navier-Stokes equations for the plane stationary
 flow of a viscous fluid can be written in the Helmholtz form as one equation
 for the  stream function $\psi$:
$$
{\partial\psi\over\partial y}
{\partial^3\psi\over\partial x^3}+
{\partial\psi\over\partial y}
{\partial^3\psi\over\partial x\partial y^2}-
{\partial\psi\over\partial x}
{\partial^3\psi\over\partial x^2\partial y}-
{\partial\psi\over\partial x}
{\partial^3\psi\over\partial  y^3}=
$$
$$
= \nu\left(
{\partial^4\psi\over\partial x^4}+2
{\partial^4\psi\over\partial x^2\partial y^2}+
{\partial^4\psi\over\partial y^4}\right),
\eqno(11)
$$
where the parameter $\nu$ is the cinematic viscous. To the  monomials of the
equation there correspond in $\R^3$ the following points
$$
Q_1=(-3,-1,2),\quad
Q_2=(-1-3,2),\quad
Q_3=Q_1,\quad
Q_4=Q_2,
$$
$$
Q_5=(-4,0,1),\quad
Q_6=(-2,-2,1),\quad
Q_7=(0,-4,1).
$$
All they ly in the plane $q_1+q_2 =-4$. Their convex hull $M$  is a
trapezoid.
To each edge of the trapezoid there corresponds a first approximation of the
equation (11). It turns out, that to the upper base
$\Gamma_1^{(1)}$ with  $D_1^{(1)} = \{Q_1=Q_3, Q_2=Q_4\}$ there correspond
the Euler equations for a nonviscous fluid. To the lower base
$\Gamma_3^{(1)}$ with  $D_3^{(1)} = \{Q_5, Q_6, Q_7\}$ there correspond the
Stokes equations for a creep flow. To the each side edge
$\Gamma_2^{(1)}$ with $D_2^{(1)} = \{Q_2=Q_4, Q_7\}$ or
$\Gamma_4^{(1)} $ with $D_4^{(1)} = \{Q_1=Q_3, Q_5\}$ there correspond the
Prandtl equations for a boundary layer.

Considering the boundary conditions from the geometry of powers point of
view, we can obtain asymptoics of the viscous flow around some particular
bodies [11].

Other examples and applications of the method for partial differential
equations see in [1, 3, 11, 19].
\bigskip

\centerline{\bf 6.~Conclusion}
\bigskip

This approach is applicable to  general systems of equations, which can be
algebraic, ordinary differential and partial differential [4, 2, 3]. It uses
methods of the geometry of power exponents developed in [1], see also [23].
We have created a program to compute the polyhedron $M$, boundary subsets
$D_k^{(d)}$ and their normal cones $U_k^{(d)}$ [12]. It was applied to a
6-dimensional problem related to the water-wade problem [13, 18]. For a
Hamiltonian system, we have given an algorithm for selecting all such its
first approximations, which are Hamiltonian systems themselfs [14].

Here we have considered an analysis of solutions: a selection of first
approximations of an equation, each of which can give only a part of a
singular solution. To obtain whole singular solution we must  match these
parts (see Introduction to  [22]). Such synthesis will be a subject of
future works.

The Newton polygon was introduced by Briot and Bouquet (1856), Newton himself
found one its edge only [26]. The Newton polyhedron was introduced in [21] for 
a
system of ODE. See its history in [1,4].

The  work was partly supported by the International Science Foundation, Grant
M3W300.

\begin{thebibliography}{99}
\bibitem{1}
    Bruno, A.D. Local Methods in Nonlinear Differential Equations.
   Springer -Verlag, Berlin, 1989.
\bibitem{2}
  Bruno, A.D. and  Soleev, A.  First approximations of algebraic equations //
  Russian Ac. Sci. Doklady. Mathem. {\bf 49} (1994).

\bibitem{3}
  Bruno, A.D.   First approximations of differential equations //
  Ibid.

\bibitem{4}
  Bruno, A.D. and  Soleev, A.  Local uniformization of branches of a space
curve, and Newton polyhedra // St. Petersburg Math. J. {\bf 3}:1~(1992)~53--82. 

\bibitem{5}
  Bruno, A.D. and  Petrovich, V.Yu.  Regularization of oscillations of a
  satellite on the very stretched orbit //Inst. Appl. Math., Preprint No.~4,
  Moscow, 1994 (Russian).

\bibitem{6}
  Bruno, A.D. and  Petrovich, V.Yu.  Computation of periodic oscillations of a
  satellite. Singular case//Ibid. No.~44, 1994 (Russian).

\bibitem{7}
  Sadov, S.Yu. Analysis of a function that defines rotational stability of
  almort symmetric satellite // Ibid. No.~84, 1994 (R).

\bibitem{8}
  Sadov, S.Yu. Coefficients of an averaged equation of satellite oscillations
  // Ibid. No.~27, 1995 (R).
\bibitem{9}
  Bruno, A.D. Singular perturbations in Hamiltonian Mechanics // Hamiltonian
  Mechanics (J. Seimenis, ed.). Plenum Press: N.Y., 1994. 43--49.

\bibitem{10}
  Bruno, A.D. Simple (double, multiple) periodic solutions of the restricted
  three-boody  problem in the Sun-Jupiter case //
  Inst. Appl. Math., Preprints Nos.~66, 67, 68,
  Moscow, 1993 (Russian).

\bibitem{11}
  Bruno, A.D. and  Vasil'ev,  M.M. Newton polyhedra and the asymptotic
analysis of the viscous fluid flow around a flat plate // Ibid. No.~44, 1995 
(R).


\bibitem{12}
 Soleev, A. and Aranson, A.B. Computation of a polyhedron and of
 normal cones of its faces. // Ibid . No~36,  1994 (R).

\bibitem{13}
 Soleev, A. and Aranson, A.B. First approximations of an invertible ODE
  system. // Ibid. No.~28,  1995 (R).


\bibitem{14}
  Bruno, A.D. and  Soleev, A. Hamiltonian truncations of a Hamiltonian system
// Ibid. No.~55, 1995 (R) = Russian Ac. Sci. Doklady. Mathem. {\bf 51} (1995). 

\bibitem{15}
  Bruno, A.D.  Bifurcation of the periodic solutions in the  symmetric  case
         of a multiple pair of  imaginary  eigenvalues // Selecta  Math.
         formerly Sovietica {\bf 12}:1 (1993) 1-12.


\bibitem{16}
  Sadov, S.Yu. On a dynamical system, arised from a
  finitedimensional aproximation of the Schr\"odinger equation // Mathematical
  Notes {\bf 56}:3 (1994).

\bibitem{16}
  Bruno A.D. and Sadov, S.Yu. On a formal integral of a nondivergent system
  // Ibid. {\bf 57}:6 (1995).

\bibitem{18}
  Bruno, A.D. and  Soleev, A.  Local analysis of a singularity of an
invertible ODE system // Institute Appl. Mathem. Preprints Nos. 40, 47, 54. 
 Moscow,
  1995 (R).


\bibitem{19}
  Bruno, A.D. The Newton polyhedron in Nonlinear Aanalysis //Ibid.
  No.~48, 1995 (R) = Vestnik of the Moscow University (1995).

\bibitem{20}
   Khovanskii, A.G. Newton polyhedra (resolution of singularities) // J.
   Sov. Math. {\bf 27} (1984) 2811--2830.

\bibitem{21}
  Bruno, A.D. The asymptotic behavior of solutions of nonlinear systems
  of differential equations // Soviet Math. Dokl. {\bf 3}(1962) 464--467.

\bibitem{22}
  Bruno, A.D. The Restricted 3-Body Problem. Walter de Gruyter, Berlin, 1994.

\bibitem{23}
  Bruno, A.D. The Geometry of power exponents // I.H.E.S., Preprint M/90/42,
  Paris, 1990.

\bibitem{24}
  Bruno, A.D. On Periodic flybys of the moon //Celestial Mechanics {\bf 24}:3
  (1981) 255--268.

\bibitem{25}
  Bruno, A.D. Local method of nonlinear analysis // Banach Center
  Publications {\bf 20}(1988) 103--120.

\bibitem{26}
  Newton, I. A treatise of the method of fluxions and infinite series, with
its application to the geometry of curve lines // The Mathematical Works of
Isaac Newton (Harry Woolf, ed.), v.1, Johnson Reprint Corp., N.Y. and London,
1964, pp.~27-137.

\end{thebibliography}

\end{document}
